\section{Reproducibility Details}

This appendix records the operational settings needed to interpret and
reproduce the completed analyses. It is intentionally more detailed than the
conference-length derivative of this manuscript.

\subsection{Cohort and Outcome Contract}

\begin{table}[!h]
\centering
\caption{Dataset roles and analysis units.}
\label{tab:data-contract}
\begin{tabularx}{\textwidth}{p{0.14\textwidth}p{0.18\textwidth}p{0.16\textwidth}X}
\toprule
Resource & Role and unit & Modalities & Outcome and boundary \\
\midrule
Coswara & Development and internal evaluation; participant is the split and
bootstrap unit & Cough, breathing, speech & Resolved positive/negative status;
unknown labels excluded from supervised analyses \\
COUGHVID & Frozen external target; recording UUID is the available unit &
Cough only & Processed \texttt{status\_SSL}; 285 positive and 8,046 negative
recordings; no target data used in source development \\
\bottomrule
\end{tabularx}
\end{table}

The external endpoint is not interchangeable with COUGHVID's expert audible-
pathology annotations or the dataset paper's cough-versus-noncough detector.
The source and target outcome fields also differ in ascertainment. External
performance therefore estimates transport to the complete COUGHVID analytic
setting, not a pure change in microphone acoustics.

\subsection{Signal and Feature Settings}

\begin{table}[!h]
\centering
\caption{Frozen conventional-audio processing settings.}
\label{tab:audio-settings}
\begin{tabular}{ll}
\toprule
Operation & Setting \\
\midrule
Waveform & Mono, 16~kHz, peak normalized \\
Low-energy trimming & 35-dB threshold \\
Active-region detector & 25-ms RMS frame, 10-ms hop, 0.25 peak ratio,
120-ms margin \\
Fixed duration & Cough 4~s; breathing 5~s; speech 5~s \\
Spectral grid & 400-sample FFT/window; 160-sample hop \\
Mel analysis & 64 bands, 20--8,000~Hz \\
Quality thresholds & Minimum 0.5~s cough or 1.0~s breath/speech; maximum
0.70 silence ratio and 0.01 clipping ratio \\
\bottomrule
\end{tabular}
\end{table}

The 10,140 numeric candidates comprised 6,377 ComParE-derived columns, 1,586
IS10-derived columns, and 2,177 project acoustic columns. The project bank
summarized MFCCs, mel energy, chroma, spectral contrast, tonal centroid,
zero-crossing rate, signal energy, spectral shape, and rhythm. These are
recording-level descriptions; they are not direct physiological measurements.

\subsection{Conventional Model Bank}

Every pipeline removed constant columns. A univariate ANOVA selector was then
fitted on training data using the model-specific retained percentage. Scaling
was used for logistic regression and SVC. SMOTE, where specified, used three
nearest neighbors and operated only after the training subset had been fixed.

\begingroup
\small
\begin{longtable}{L{0.22\textwidth}L{0.12\textwidth}L{0.56\textwidth}}
\caption{Principal conventional classifier settings.}
\label{tab:model-settings}\\
\toprule
Model & ANOVA features & Main settings \\
\midrule
\endfirsthead
\toprule
Model & ANOVA features & Main settings \\
\midrule
\endhead
L2 logistic & 80\% & $C=1$; class-balanced; maximum 4,000 iterations \\
RBF SVC & 60\% & $C=2$; scale-based $\gamma$; class-balanced; three-fold
sigmoid probability calibration \\
Extra Trees & 100\% & 700 trees and class balancing without SMOTE, or 800
trees with SMOTE; minimum leaf size 2 \\
Random Forest & 80\% & 700 trees with balanced subsampling, or 800 trees with
SMOTE; minimum leaf size 2 \\
XGBoost+SMOTE & 80\% & 800 trees; depth 3; learning rate 0.025; row and column
subsampling 0.85; minimum child weight 2; L2 penalty 2 \\
LightGBM+SMOTE & 80\% & 900 trees; 31 leaves; learning rate 0.025; minimum 20
child samples; row and column subsampling 0.85; L2 penalty 2 \\
CatBoost+SMOTE & 80\% & 900 iterations; depth 5; learning rate 0.025; log-loss
training objective \\
\bottomrule
\end{longtable}
\endgroup

The development bank also included SMOTE and non-SMOTE variants and a
validation-only Optuna search over logistic regression, Extra Trees, and
Random Forest. Model identity, modality combination, fusion method, and
threshold were selected from validation evidence. External labels did not
participate in this process.

\subsection{Deep Representation Settings}

WavLM Base Plus operated on 3-s windows with 50\% overlap and at most four
segments per recording. One augmented training copy was permitted. Training
used batch size 8, gradient accumulation 4, learning rates $2\times10^{-5}$
for the encoder and $10^{-4}$ for the prediction head, four unfrozen top
transformer layers, eight maximum epochs, and early stopping patience 3. The
CNN--BiGRU used log-mel spectrograms, batch size 16, learning rate $10^{-3}$,
and eight epochs. Both analyses selected checkpoints and thresholds using
Coswara validation data before frozen COUGHVID transfer.

\subsection{Interpretive Boundaries}

The following distinctions are part of the analysis contract:

\begin{itemize}
    \item the 0.897 row is internal cough--speech fusion, whereas COUGHVID tests
    only a cough branch;
    \item the early-to-late selected system is not the same model and modality
    combination as the internal selected fusion;
    \item independent source and target samples support independent bootstrap
    differences, not paired DeLong testing;
    \item metadata predictiveness is evidence of collection-context
    association, not proof that the audio model encoded the same variables;
    and
    \item post-hoc mechanism, subgroup, and utility analyses are exploratory.
\end{itemize}
